% Mathématiques 4e — cours V3
% Statistiques et médiane

\section{Objectifs}

\begin{itemize}
\item Recueillir et organiser une série de données.
\item Calculer et interpréter des effectifs et des fréquences.
\item Calculer une moyenne simple ou pondérée par des effectifs.
\item Ordonner une série.
\item Déterminer et interpréter une médiane.
\item Comparer moyenne et médiane.
\item Lire et produire tableaux, diagrammes et graphiques.
\item Choisir un indicateur ou une représentation adaptée à une question.
\item Exercer un regard critique sur des données et leurs représentations.
\end{itemize}

\section{1. Une série statistique}

Une série statistique est un ensemble de données recueillies sur une population.

Par exemple, on peut relever le temps de trajet de chaque élève d'une classe :
\[
12,\ 18,\ 9,\ 15,\ 11,\ 35,\ldots
\]

Avant tout calcul, il faut connaître :
\begin{itemize}
\item la population ;
\item le caractère étudié ;
\item l'unité ;
\item l'effectif total.
\end{itemize}

\section{2. Effectif}

L'effectif d'une valeur est le nombre de fois où cette valeur apparaît.

Pour la série :
\[
2,\ 3,\ 3,\ 4,\ 4,\ 4,\ 7,
\]
l'effectif de :
\[
4
\]
est :
\[
3.
\]

L'effectif total vaut :
\[
7.
\]

\section{3. Fréquence}

La fréquence d'une valeur ou catégorie est :
\[
\boxed{
f=\dfrac{\text{effectif de la valeur}}{\text{effectif total}}
}.
\]

Dans l'exemple précédent :
\[
f(4)=\dfrac37.
\]

On peut aussi l'exprimer en écriture décimale ou en pourcentage.

\section{4. Tableau d'effectifs}

Une série peut être regroupée.

\begin{tabular}{c|c|c}
Valeur & Effectif & Fréquence \\
$10$ & $3$ & $15\%$ \\
$11$ & $5$ & $25\%$ \\
$12$ & $8$ & $40\%$ \\
$13$ & $4$ & $20\%$ \\
\end{tabular}

L'effectif total est :
\[
3+5+8+4=20.
\]

La somme des fréquences vaut :
\[
100\%.
\]

\section{5. Moyenne simple}

Pour $n$ valeurs :
\[
x_1,\ldots,x_n,
\]
la moyenne est :
\[
\boxed{
\overline{x}=\dfrac{x_1+\cdots+x_n}{n}
}.
\]

Pour :
\[
8,\ 10,\ 13,\ 9,
\]
la moyenne vaut :
\[
\dfrac{8+10+13+9}{4}=10.
\]

\section{6. Moyenne avec effectifs}

Lorsque les valeurs sont regroupées par effectifs :
\[
\boxed{
\overline{x}
=
\dfrac{x_1n_1+x_2n_2+\cdots+x_kn_k}{n_1+\cdots+n_k}
}.
\]

\begin{exemple}
Les notes $8$, $10$ et $14$ apparaissent respectivement $2$, $3$ et $1$ fois.

La moyenne vaut :
\[
\dfrac{8\times2+10\times3+14\times1}{6}
=
\dfrac{60}{6}
=
10.
\]
\end{exemple}

\section{7. Ordonner une série}

La médiane exige d'abord de ranger les valeurs.

Série brute :
\[
12,\ 5,\ 8,\ 14,\ 7.
\]

Série ordonnée :
\[
5,\ 7,\ 8,\ 12,\ 14.
\]

La position des valeurs devient alors exploitable.

\section{8. Médiane : idée}

\begin{definition}
Une médiane est une valeur qui partage une série ordonnée en deux groupes contenant au moins la moitié des données de chaque côté.
\end{definition}

Pour un effectif impair, la valeur centrale fournit immédiatement la médiane.

\coursefigure{/images/mathematiques/4eme/statistiques-mediane.svg}{Série de neuf valeurs ordonnées dont la cinquième est mise en évidence.}{Avec neuf valeurs, la cinquième laisse quatre valeurs de chaque côté : c'est la médiane.}

\section{9. Effectif impair}

Considérons :
\[
4,\ 6,\ 7,\ 8,\ 10,\ 12,\ 14,\ 17,\ 21.
\]

Il y a :
\[
9
\]
valeurs.

La valeur centrale est la cinquième :
\[
\boxed{10}.
\]

Il reste quatre valeurs avant et quatre valeurs après.

\section{10. Effectif pair}

Considérons :
\[
3,\ 5,\ 8,\ 9,\ 12,\ 15.
\]

Il y a :
\[
6
\]
valeurs.

Les deux valeurs centrales sont :
\[
8
\]
et :
\[
9.
\]

Au collège, on choisit couramment comme médiane la moyenne de ces deux valeurs :
\[
\dfrac{8+9}{2}=8{,}5.
\]

Donc :
\[
\boxed{8{,}5}.
\]

\section{11. Interpréter la médiane}

Si la médiane d'un temps de trajet vaut :
\[
18\ \text{min},
\]
cela signifie qu'au moins la moitié des trajets ont une durée inférieure ou égale à :
\[
18\ \text{min},
\]
et qu'au moins la moitié ont une durée supérieure ou égale à :
\[
18\ \text{min}.
\]

La médiane ne signifie pas que tous les trajets sont proches de $18$ minutes.

\section{12. Moyenne et médiane}

La moyenne utilise les valeurs numériques de toutes les données.

La médiane utilise principalement leur ordre.

\coursefigure{/images/mathematiques/4eme/statistiques-moyenne-mediane.svg}{Deux distributions illustrant l'influence d'une valeur extrême sur la moyenne.}{Une valeur très éloignée peut déplacer fortement la moyenne sans nécessairement déplacer la médiane autant.}

Ces deux indicateurs peuvent donc apporter des informations différentes.

\section{13. Valeur extrême}

Considérons :
\[
10,\ 11,\ 11,\ 12,\ 56.
\]

La médiane vaut :
\[
11.
\]

La moyenne vaut :
\[
\dfrac{10+11+11+12+56}{5}
=
20.
\]

La valeur $56$ augmente fortement la moyenne.

La médiane reste proche des quatre premières valeurs.

\section{14. Comparer deux séries}

Deux classes peuvent avoir la même moyenne et des médianes différentes.

Deux séries peuvent aussi avoir la même médiane mais des moyennes différentes.

Un seul indicateur ne décrit donc pas complètement une série.

Il faut revenir à la question posée.

\section{15. Diagramme en barres}

Un diagramme en barres permet de comparer des valeurs ou catégories.

La hauteur des barres traduit l'effectif ou la fréquence.

On contrôle :
\begin{itemize}
\item le titre ;
\item l'échelle ;
\item les unités ;
\item l'origine de l'axe.
\end{itemize}

\section{16. Diagramme circulaire}

Le disque entier représente :
\[
100\%.
\]

Pour une fréquence :
\[
35\%,
\]
l'angle du secteur vaut :
\[
0{,}35\times360^\circ=126^\circ.
\]

Un diagramme circulaire est adapté à une répartition d'un total.

\section{17. Histogramme ou barres ?}

Dans une situation de données regroupées en classes d'intervalles, on peut rencontrer des représentations où les rectangles sont jointifs.

Il ne faut pas confondre cette situation avec un simple diagramme de catégories indépendantes.

En 4e, l'essentiel est surtout de lire correctement la représentation proposée.

\section{18. Graphique d'évolution}

Une courbe peut représenter une grandeur au cours du temps.

La moyenne ou la médiane d'une série n'est pas nécessairement visible directement sur cette courbe.

Chaque outil répond à une question différente :
\[
\text{évolution}\neq\text{position centrale}.
\]

\section{19. Données manquantes}

Si des données sont absentes, il faut le signaler.

Une moyenne calculée sur :
\[
28
\]
réponses n'est pas automatiquement représentative d'une classe de :
\[
32
\]
élèves si les quatre non-réponses ne sont pas aléatoires.

Même au collège, la qualité des données compte.

\section{20. Choisir un indicateur}

\begin{methode}
\begin{itemize}
\item Pour partager équitablement un total : la moyenne est naturelle.
\item Pour chercher une valeur centrale résistante à une valeur extrême : la médiane peut être plus pertinente.
\item Pour connaître une proportion : utiliser une fréquence.
\item Pour comparer des catégories : utiliser un tableau ou un diagramme adapté.
\end{itemize}
\end{methode}

\section{21. Exemple développé}

Les durées en minutes sont :
\[
8,\ 10,\ 11,\ 11,\ 13,\ 15,\ 16,\ 40.
\]

L'effectif est pair :
\[
8.
\]

Les valeurs centrales sont :
\[
11
\]
et :
\[
13.
\]

La médiane vaut :
\[
\dfrac{11+13}{2}=12.
\]

La moyenne vaut :
\[
\dfrac{124}{8}=15{,}5.
\]

La valeur $40$ tire la moyenne vers le haut.

\section{22. Tableur}

Un tableur permet :
\begin{itemize}
\item de trier les données ;
\item de calculer une moyenne ;
\item de compter des effectifs ;
\item de calculer des fréquences ;
\item de produire un graphique.
\end{itemize}

Mais l'élève doit encore interpréter le résultat et choisir l'indicateur pertinent.

\section{23. Méthode de médiane}

\begin{methode}
\begin{enumerate}
\item ranger toutes les valeurs dans l'ordre croissant ;
\item compter l'effectif total ;
\item si l'effectif est impair, repérer la valeur centrale ;
\item s'il est pair, repérer les deux valeurs centrales et appliquer la convention utilisée ;
\item interpréter le résultat dans le contexte.
\end{enumerate}
\end{methode}

\section{24. Erreurs fréquentes}

\begin{itemize}
\item Chercher la médiane sans ordonner la série.
\item Confondre médiane et moyenne.
\item Prendre la moitié de la plus grande valeur.
\item Oublier les effectifs dans un calcul de moyenne pondérée.
\item Interpréter une fréquence comme un effectif.
\item Lire un graphique sans vérifier l'échelle.
\item Présenter une moyenne sans unité alors que les données en ont une.
\end{itemize}

\section{À retenir}

\begin{aretenir}
La fréquence vaut :
\[
\boxed{f=\dfrac{\text{effectif}}{\text{effectif total}}}.
\]

La moyenne simple vaut :
\[
\boxed{\overline{x}=\dfrac{\text{somme des valeurs}}{\text{effectif total}}}.
\]

La médiane partage une série ordonnée en deux moitiés.

Moyenne et médiane ne décrivent pas la même propriété d'une série : il faut choisir l'indicateur selon la question et interpréter les résultats dans leur contexte.
\end{aretenir}
